\section{Dataset and preprocessing}
\label{sec:data}

\subsection{Cohort construction and clinical variables}
We retrospectively analysed longitudinal breast dynamic contrast-enhanced MRI
(DCE-MRI) from the public I-SPY2 and ACRIN-6698 imaging
resources~\cite{tciaispy2_2022,tciaacrin6698_2021}. Source imaging identifiers
were linked to research identifiers using the accompanying official identity
mapping, and outcomes and baseline clinical variables were obtained from the
official I-SPY2 Imaging Cohort~1 clinical spreadsheet. The final cohort
comprised 976 patients with 3,641 usable examinations. Native imaging
identifiers belonged to the ISPY2 and ACRIN-6698 namespaces for 712 and 264
patients, respectively; these identifiers were mapped to the same clinical
release and do not define independent development and external validation
cohorts. The official collection pages identify the source imaging and clinical
release as CC BY 4.0. \todo{Confirm terms for redistributing derived assets and
document the applicable ethics approval or waiver and original consent arrangements.}

Eligibility required an actual baseline examination, an available binary pCR
label, usable baseline registered image assets, and a nonempty baseline
localisation produced by the existing MAMA-MIA segmentation workflow. Of 980
clinically matched source candidates, three lacked a baseline examination and
one had an empty or unavailable predicted baseline localisation. Their
exclusion yielded 976 patients and 3,642 candidate examinations. One
post-treatment examination was subsequently marked unavailable because its
saved within-examination phase transforms displaced the postcontrast images
outside the prescribed baseline crop. The patient's earlier examinations were
retained, leaving 3,641 usable examinations. Missing follow-up examinations did
not cause patient exclusion. No additional admission thresholds based on
normalised cross-correlation, field-of-view ratios, functional tumour volume,
or the availability of an accepted adjacent longitudinal pair were applied.

The endpoint was the released binary pCR field: 313 patients were labelled
positive and 663 negative. \todo{Verify and state the release-specific
pathological definition of pCR, including breast and nodal criteria and the
handling of in-situ disease; do not infer this definition from the binary
column alone.} Baseline covariates were age at screening, hormone-receptor
(HR) status, HER2 status, and the released binary MammaPrint indicator. Age was
available for 973 patients; all three binary covariates were available for the
entire cohort. The official assigned treatment arm, represented by 13
categories, was also supplied to the model. The current experiments use an
\emph{assigned-arm scenario}: the assigned arm and available baseline
covariates are assumed known at baseline. The spreadsheet does not establish
prospective allocation timestamps, doses, cycle-level exposure, or interval
treatment schedules, and the model does not reconstruct these quantities.

\begin{table*}[t]
\centering
\caption{Characteristics of the analysed cohort. Development statistics count
each of the 801 development patients once, rather than pooling overlapping
training folds. Age is reported as mean $\pm$ sample standard deviation among
patients with observed age; other entries are counts.}
\label{tab:cohort}
\small
\begin{tabular}{lrrr}
\toprule
Characteristic & All & Development & Holdout \\
\midrule
Patients & 976 & 801 & 175 \\
pCR positive & 313 & 257 & 56 \\
pCR negative & 663 & 544 & 119 \\
Age (years) & $48.7\pm10.5$ & $48.7\pm10.6$ & $48.9\pm10.2$ \\
Missing age & 3 & 3 & 0 \\
HR positive & 533 & 434 & 99 \\
HER2 positive & 240 & 195 & 45 \\
Available T0 & 976 & 801 & 175 \\
Available T1 & 935 & 760 & 175 \\
Available T2 & 875 & 700 & 175 \\
Available T3 & 855 & 692 & 163 \\
\bottomrule
\end{tabular}
\end{table*}

Examinations were indexed by the canonical study stages T0--T3 rather than
continuous elapsed time. Availability at these stages was 976, 935, 875, and
855 patients, respectively (Table~\ref{tab:cohort}); 803 patients had all four
examinations and 861 had T0--T2. The nominal source protocol acquired T0 before
randomization, T1 after three weeks of the first regimen (paclitaxel with or
without an experimental agent), T2 between paclitaxel and anthracycline
regimens, and T3 after four cycles of anthracycline treatment and before
surgery~\citep{tciaispy2_2022}. Because T0 MRI preceded randomization, supplying
the eventual assigned arm at T0 is a retrospective conditional scenario,
rather than a prospectively available predictor before treatment allocation.
The source protocol used 1.5 or 3.0 Tesla scanners with dedicated breast coils
and required the same scanner configuration within a patient. These are
nominal protocol specifications, not a patient-level audit of timing or
scanner compliance. Twenty-three patients had T0 only and remained
eligible for pCR supervision. Their unavailable future MRI targets were
masked, and no synthetic follow-up images were substituted. Model inputs and
evaluation landmarks were restricted to T0, T1, and T2; a real T3 examination,
when available in the training partition, could provide a future-state target.
\todo{Report the actual examination intervals, acquisition centres, scanner
distribution, and sequence-parameter ranges for the final analysed cohort.}

\subsection{Patient partitioning and evaluation scope}
The cohort was divided into 801 development patients and a fixed internal
holdout of 175 patients, including 56 pCR-positive patients. All holdout
patients had actual T0, T1, and T2 examinations, allowing matched longitudinal
evaluation; 163 also had T3. The holdout retained 139 patients with complete
T0--T2 observations from the historical benchmark and added 36 eligible
patients. Additional patients were selected using outcome quotas and a fixed
baseline-clinical distribution balancing rule with split seed 2026, without
using model predictions or performance. This fixed cohort-construction seed
is distinct from the training-seed comparison: the split is unchanged for
training seeds 42, 44, 71, and 2026. Development folds were stratified by
pCR and examination-availability pattern; rare joint strata were combined
within the same pCR class. Each development patient appeared in validation
exactly once across the five folds. Fold~1 contained 640 training and 161
validation patients; the remaining folds each contained 641 training and 160
validation patients. Training, validation, and holdout identities were
disjoint within each fold, and the same holdout was used for every fold
model at each training seed.

This partition is an internal downstream benchmark. In particular, 29 of the
36 added holdout patients had belonged to the development population of an
earlier experiment, so earlier trained models cannot be treated as held out
on the entire new set. The present downstream models were trained using the
new fold definitions. Patient independence of the supplied codec and
segmentation pretraining checkpoints relative to the final holdout has not
yet been established; that provenance must be audited before stronger claims
of end-to-end independence. The complete-examination holdout eligibility also
restricts inference about patients with missing early follow-up examinations.

\subsection{Registration, fixed localisation, and normalisation}
Within each examination, postcontrast phases were registered to the
precontrast reference. Follow-up examinations were aligned to the patient's
T0 reference using the saved longitudinal transforms, retaining a rigid
fallback when the final deformable candidate was not usable. Baseline
localisation used the existing MAMA-MIA nnU-Net prediction on the native
first-postcontrast examination~\cite{mamamia2025}. The centre of the largest
26-connected predicted component defined a fixed baseline crop, which was
reused at every actual follow-up. Follow-up segmentation and follow-up-driven
recentering were not used.

Each examination was represented by three channels: the precontrast phase,
the first postcontrast phase, and the late phase explicitly specified in the
examination metadata. The late channel was not selected by assuming a fixed
phase number or taking the last available phase. Images were resampled to
$32\times128\times128$ voxels with spacing
$2.0\times0.7032\times0.7032\,\mathrm{mm}$ in the slice and two in-plane
dimensions, approximately covering $64\times90\times90\,\mathrm{mm}$.
Linear interpolation was used for MRI intensities and nearest-neighbour
interpolation for masks and acquisition support, with zero padding outside
the acquired field of view. The localisation mask determined the crop; it was
not a contemporaneous ground-truth tumour segmentation or a future
observation supplied to the classifier.

Native metadata, baseline mask headers, registered reference grids, and saved
transforms were maintained as matched versions. Explicit array-index
mappings, rather than placeholder NIfTI affines, defined sampling positions.
The exported arrays were checked for finite values, consistent coordinate
mapping, and numerical agreement after compressed storage and readback.
These checks establish data and geometry consistency but do not establish
anatomical registration or segmentation accuracy. The fixed crop partly
truncated the largest predicted baseline component for 85 patients; these
patients were retained and the limitation was recorded.

Raw cropped intensities were stored once and shared across folds. For each
fold, a single intensity mean and population standard deviation were fitted
to nonzero precontrast foreground voxels from every unique actual examination
of that fold's training patients. The same scalar transform was applied once
to the foreground of all three channels, stages, and partitions; background
and padding remained zero. Validation and holdout images were excluded from
fitting. Age normalisation was likewise fitted to observed training-patient
values; HR, HER2, and MammaPrint retained their official binary coding, with
missingness represented explicitly.
